\section{Monte Carlo design}\label{sec:monte_carlo}

\paragraph{Streams.} Each random source has its own PCG64 generator derived
from a single master \texttt{SeedSequence}: \texttt{price}, \texttt{rate},
\texttt{credit}, \texttt{lgd}, \texttt{qmc}. This guarantees
\emph{bit-identical} reproducibility across runs while keeping the streams
independent, which is the requirement for the variance-reduction
techniques to behave as expected (see
\texttt{tests/unit/test\_engine.py}).

\paragraph{Antithetic — formal variance reduction.}
Let $f$ be a measurable function of an independent Gaussian sample
$Z_1, \ldots, Z_n$. For any $i$ the antithetic pair $(Z_i, -Z_i)$ has
perfect negative correlation $\rho = -1$. For a linear functional
$f(z) = a + b\,z$ the pair average $\bar f_i = (f(Z_i) + f(-Z_i))/2 = a$
is deterministic and the variance vanishes; for any monotone $f$ the
Hoeffding lemma gives $\Cov(f(Z), f(-Z)) \le 0$, so $\Var(\bar f) \le
\Var(f(Z))/2$. Glasserman (2003, Section 4.2) formalises the general
result: for any $f$ that is monotone in each component of $z$,
\begin{equation*}
\Var\!\left(\frac{1}{n}\sum_i \bar f_i\right) \le \Var\!\left(\frac{1}{n}\sum_i f(Z_i)\right),
\end{equation*}
so antithetic never hurts on monotone estimands and strictly helps on
linear ones. The senior tranche expected loss is approximately linear
in the common factor below the attach point, which is why
\cref{fig:mc-convergence} shows the antithetic standard error roughly
five times below the plain MC at every sample size.

\paragraph{Control variate from the Vasicek (2002) LPF.}
The large-portfolio limit \cref{eq:vasicek-lpf} delivers a closed-form
expected loss $\mathbb{L}^{\mathrm{LPF}} = \E[L \mid \text{LPF}]$ that is
highly correlated with the MC estimator $\hat L_n$ on finite portfolios.
The control-variate-adjusted estimator
\begin{equation*}
\hat L_n^{\mathrm{cv}} = \hat L_n - \hat\beta \bigl(\mathbb{L}^{\mathrm{LPF}}_{\text{MC}} - \mathbb{L}^{\mathrm{LPF}}\bigr),
\quad \hat\beta = \frac{\widehat{\Cov}(L, L^{\mathrm{LPF}})}{\widehat{\Var}(L^{\mathrm{LPF}})}
\end{equation*}
has variance reduced by a factor $1 - \rho_{L, L^\mathrm{LPF}}^2$ relative
to $\hat L_n$. In our 70-obligor portfolio the empirical correlation is
above 0.95, so the variance reduction is comparable to the antithetic
gain. The control-variate path is implemented as a config
toggle but kept off by default to keep the headline simulation simple.

\paragraph{Quasi-Monte Carlo.}
For the copula common factor we provide a Sobol-based QMC alternative
\autocite{Owen1997}. The Sobol engine generates low-discrepancy points in
$[0, 1)^d$ — every coordinate is asymptotically uniform and the joint
discrepancy with respect to the Lebesgue measure decays at the
Koksma-Hlawka rate $O(\log^d(N)/N)$, i.e.\ $O(N^{-1})$ up to
logarithmic factors once $N$ grows past the dimension. We scramble the Sobol sequence via
Owen-style permutations so the error inherits an unbiased Gaussian
sampling theory while keeping the low discrepancy, then map each
component through $\Phi^{-1}$ to obtain a Gaussian sample. The
implementation lives in \texttt{simulation/qmc.py}; activating it
requires setting \texttt{use\_qmc=True} in \texttt{SimulationConfig}. The
empirical gain is largest when the integrand is smooth in the common
factor — for the tail-heavy senior-loss estimator the practical gain is
modest, as confirmed in \cref{fig:mc-convergence}.

\paragraph{Empirical convergence.} \Cref{fig:mc-convergence} reports the
realised standard error of the senior tranche expected PV over $K = 12$
independent simulations at each sample size $N \in \{256, 512, 1024,
2048, 4096, 8192\}$. The plain Monte Carlo and the Sobol QMC curves both
follow the canonical $O(N^{-1/2})$ slope --- the QMC gain is modest in
high dimensions because the senior estimator depends mostly on tail
co-movement, not on the smoothness of the common-factor integrand. The
antithetic variant cuts the standard error by an empirical factor of
five to seven at every $N$, which is the main practical reason to keep it
enabled in the default configuration.

\begin{figure}[ht]
  \centering
  \includegraphics[width=0.85\linewidth]{fig_mc_convergence.pdf}
  \caption{Empirical standard error of the senior tranche expected PV as
  a function of the number of Monte Carlo paths. Reference slopes
  $O(N^{-1/2})$ (dashed) and $O(N^{-1})$ (dotted) anchored at the
  leftmost point.}
  \label{fig:mc-convergence}
\end{figure}
